\chapter{Quantities and measurements}

\section{Quantities and units}
\bf{Physical quantities} can only be defined in terms of other physical quantities, usually base quantities.
\begin{example}{}{derived-quantities}
	The definition of \it{velocity} is `the rate of change in \it{displacement}
	with respect to \it{time}', and not `the change in displacement per second'.
\end{example}

\begin{table}[htpb]
	\centering
	\begin{tabular}{c c}
		\toprule
		base quantity             & base unit        \\
		\midrule
		mass                      & \unit{\kilogram} \\
		length                    & \unit{\metre}    \\
		time                      & \unit{\second}   \\
		electric current          & \unit{\ampere}   \\
		thermodynamic temperature & \unit{\kelvin}   \\
		amount of substance       & \unit{\mol}      \\
		luminous intensity        & \unit{\candela}  \\
		\bottomrule
	\end{tabular}
	\label{tab:base-quantities}
\end{table}

\begin{definition}{Homogeneity}{homogeneous}
	A physical equation is \bf{homogeneous} if all terms have the same units, when
	expressed in SI base units.
\end{definition}
\Cref{def:homogeneous} is necessary for a physical equation to be valid.

\begin{table}[htpb]
	\centering
	\begin{tabular}{c c}
		SI prefix & factor     \\
		\toprule
		pico      & \num{e-12} \\
		nano      & \num{e-9}  \\
		micro     & \num{e-6}  \\
		milli     & \num{e-3}  \\
		centi     & \num{e-2}  \\
		deci      & \num{e-1}  \\
		kilo      & \num{e3}   \\
		mega      & \num{e6}   \\
		giga      & \num{e9}   \\
		tera      & \num{e12}  \\
		\bottomrule
	\end{tabular}
\end{table}

\section{Error and uncertainty}
A measurement \(X\) is reported as
\begin{equation*}
	X = x \pm \Delta x,
\end{equation*}
where \(x\) is the \bf{measured value} and \(\Delta x\) is the \bf{absolute uncertainty}.
\(\lf{\Delta x}{x}\) is the \bf{fractional uncertainty}.

\(\Delta x\) must be expressed to 1 significant figure, after which \(x\) must be rounded to the same decimal place as \(\Delta x\).
\begin{table}[htpb]
	\centering
	\begin{tabular}{c c}
		\toprule
		operation         & uncertainty                                                                                 \\
		\midrule
		\(z = x\pm y\)    & \(\Delta z = \Delta x + \Delta y\)                                                          \\
		\(k = x^ay^bz^c\) & \(\f{\Delta k}{k} = \ab|a|\f{\Delta x}{x} + \ab|b|\f{\Delta y}{y} + \ab|c|\f{\Delta z}{z}\) \\
		\(y = f\ab(x)\)   & \(\Delta y = \f{1}{2}\ab(\max y - \min y)\)                                                 \\
		\bottomrule
	\end{tabular}
	\label{eq:uncertainty}
\end{table}

\begin{definition}{Systematic error}{systematic-error}
	\bf{Systematic errors} are deviations of the mean value from the true value of a measurement, with \it{fixed direction}
	and \it{similar magnitude}.

	They can be reduced by improving the experimental design.
\end{definition}
\begin{definition}{Random error}{random-error}
	\bf{Random errors} are deviations of the mean value from the true value of a measurement
	that \it{vary in magnitude and direction}.

	They can be reduced by finding the mean of repeated measurements.
\end{definition}
\begin{definition}{Precision}{precision}
	The degree of consistency and agreement among independent measurements of the same quantity.
\end{definition}
\begin{definition}{Accuracy}{accuracy}
	The closeness of agreement between a measured value and the true value of a quantity.
\end{definition}